\section{Related Work}
\label{sec:related}

We briefly discuss three directions most relevant to this work: unified dense grounded MLLMs, visual token compression, and learnable budgeted selection.
Our setting differs from standard VQA-style compression because token reduction must preserve segmentation quality, temporal grounding, and multimodal conversation within one model.

\paragraph{Dense grounded multimodal understanding.}
Dense grounding extends multimodal understanding from text answers to pixel-level outputs.
Recent works explore this direction from image grounding, region captioning, and referring segmentation perspectives, including Osprey~\cite{yuan2024osprey}, LISA~\cite{lai2023lisa}, GLaMM~\cite{rasheed2024glamm}, OMG-LLaVA~\cite{zhang2024omgllava}, and VISA~\cite{yan2024visa}.
At the same time, unified MLLMs such as LLaVA~\cite{liu2023visual} and LLaVA-OneVision~\cite{li2024llavaonevision} push toward broader task coverage across image and video inputs.
Sa2VA~\cite{yuan2025sa2va} combines these trends by integrating SAM2~\cite{ravi2024sam2} with an MLLM and exposing a shared image-video dense grounding interface.
This unified design is what makes token reduction more delicate in our setting: the compressed visual stream must still preserve spatial detail, temporal continuity, and prompt-conditioned evidence for both conversation and mask decoding.

\paragraph{Visual token compression in MLLMs.}
Visual token compression aims to reduce the number of tokens processed by the LLM while preserving task performance.
Earlier approaches often rely on heuristic pruning, fixed downsampling, token merging, or architecture-specific compression rules~\cite{yang2024visionzip,zhang2024sparsevlm,zhu2024focusllava,yao2024deco}.
These methods reduce compute effectively, but are mostly validated on question-answering-oriented MLLMs.
More recent approaches move toward adaptive compression and inference-compatible pruning~\cite{yang2025topv,zhu2025visionselector}.
VisionSelector~\cite{zhu2025visionselector} is closest to our goal because it couples a learnable scorer with differentiable Top-$k$ training and hard inference-time pruning.
Our difference is the target model: instead of compressing a general VQA-style MLLM, we study whether learnable selection remains reliable inside Sa2VA's segmentation-sensitive unified visual stream.

\paragraph{Learnable selection for unified dense grounding.}
The key distinction of our work is therefore not a new selector algorithm, but a new transfer setting.
Sa2VA must jointly handle segmentation and conversation, and its visual input mixes ordinary image/video tokens with visual-prompt-conditioned tokens.
This makes token selection more tightly coupled to mask quality and temporal grounding than in conventional multimodal QA.
Accordingly, we do not present LIS, DiffTopK, or CAS as newly invented here.
Our contribution is to integrate learnable token selection into Sa2VA's unified visual stream in a compatible way, and to verify experimentally when this transfer improves dense grounding under explicit token budgets.
